\chapter{Cholecystitis}
\index{Cholecystitis}
\label{sec:Cholecystitis}

\section{Overview}
Acute calculous cholecystitis results from persistent impaction of a gallstone in the cystic duct, causing bile stasis, mucosal reduced blood flow, chemical irritation, and secondary bacterial superinfection (Escherichia coli, Klebsiella, Enterococcus). Acalculous cholecystitis occurs in critically ill patients (burns, trauma, major surgery, prolonged parenteral nutrition).
\section{Causes}
\section{Risk Factors}

\begin{itemize}[noitemsep]
  \item Genetic predisposition and family history
  \item Advancing age
  \item Lifestyle factors (diet, exercise, smoking, alcohol)
  \item Environmental exposures
  \item Underlying medical conditions
  \item Medication use
\end{itemize}
\section{Signs and Symptoms}

\subsection{Common symptoms}
\begin{itemize}[noitemsep]
  \item You may experience severe steady right upper quadrant (RUQ) or epigastric pain radiating to the right scapula/shoulder, fever, leukocytosis, and a positive Murphy sign (inspiratory arrest on deep palpation of the RUQ). Diagnosis is supported by abdominal ultrasonography showing gallbladder wall thickening ($>4$ mm), pericholecystic fluid, and sonographic Murphy sign
  \item HIDA scan confirms cystic duct obstruction if ultrasound is equivocal.
  \item Pain or discomfort in the affected area
  \end{itemize}

\subsection{Emergency warning signs}
\begin{itemize}[noitemsep]
  \item Severe or worsening symptoms of cholecystitis require immediate medical evaluation
\end{itemize}
\section{Progression and Stages}
Cholecystitis is inflammation of the gallbladder wall, classified as acute calculous (90\% of cases), acalculous (10\%), or chronic. The condition may progress through different stages of severity over time. Early stages often involve mild or subtle symptoms that may be overlooked. Moderate stages involve more noticeable symptoms affecting daily function. Advanced stages may involve significant impairment and complications.
\section{Complications}
Complications include gallbladder gangrene, perforation, and emphysematous cholecystitis.

\begin{itemize}[noitemsep]
  \item Disease progression leading to worsening symptoms
  \item Secondary infections or injuries
  \item Side effects of treatment
  \item Reduced quality of life
  \item Emotional and psychological impact
\end{itemize}
\section{Diagnosis}
Comprehensive medical history and physical examination Laboratory tests to assess organ function and rule out other conditions Imaging studies to visualize affected structures Specialized testing depending on the affected system Consultation with relevant specialists Monitoring of disease progression over time

\begin{itemize}[noitemsep]
  \item Comprehensive medical history and physical examination
  \item Laboratory tests to assess organ function
  \item Imaging studies to visualize affected structures
  \item Specialized testing depending on the affected system
  \item Consultation with relevant specialists
\end{itemize}
\section{Prevention}
\begin{itemize}[noitemsep]
  \item Maintain a healthy lifestyle with balanced diet and exercise
  \item Avoid known risk factors
  \item Attend regular medical check-ups for early detection
  \item Follow recommended screening guidelines
\end{itemize}
\section{Treatment Options}
Treatment options include NPO status, IV fluids, IV broad-spectrum antibiotics, and early using small incisions and a camera cholecystectomy (within 24--72 hours). Medications to manage symptoms and address underlying mechanisms Lifestyle modifications and self-care measures Physical therapy and rehabilitation Surgical intervention when indicated Regular monitoring and follow-up care Patient education and support services

\begin{itemize}[noitemsep]
  \item Medications to manage symptoms and address underlying mechanisms
  \item Lifestyle modifications and self-care measures
  \item Physical therapy and rehabilitation
  \item Surgical intervention when indicated
  \item Regular monitoring and follow-up care
\end{itemize}
\section{Living with It}
\begin{itemize}[noitemsep]
  \item Follow your treatment plan as prescribed
  \item Maintain regular follow-up with your healthcare provider
  \item Make necessary lifestyle adjustments
  \item Seek support from family, friends, or support groups
  \item Stay informed about your condition
\end{itemize}
\section{Outlook}
The outlook varies depending on the specific condition, severity at diagnosis, and response to treatment. Early intervention generally leads to better outcomes. Many conditions can be effectively managed with appropriate treatment. Regular follow-up care is important for monitoring and treatment adjustment.
